\documentclass[a4paper,11pt,reqno]{amsart}

\usepackage[T1]{fontenc}
\IfFileExists{lmodern.sty}
  {\usepackage{lmodern}}
  {\PassOptionsToPackage{expansion=false}{microtype}}
\usepackage{microtype}
\usepackage{mathtools,amssymb,amsthm}
\usepackage{xcolor}
\usepackage[a4paper,left=21mm,right=21mm,top=17mm,bottom=19mm]{geometry}
\usepackage[colorlinks=true,allcolors=blue!55!black]{hyperref}

\newtheorem{maintheorem}{Theorem}
\renewcommand{\themaintheorem}{\Alph{maintheorem}}
\newtheorem{theorem}{Theorem}
\newtheorem{lemma}[theorem]{Lemma}
\newtheorem{corollary}[theorem]{Corollary}

\colorlet{theoremframe}{blue!55!black}
\colorlet{theoremfill}{blue!4}
\colorlet{corollaryframe}{green!40!black}
\colorlet{corollaryfill}{green!5}
\colorlet{aiframe}{red!70!black}
\colorlet{aifill}{yellow!18}

\newcommand{\leanrepo}{https://github.com/TimmPeterson/lie-algebras/blob/da1384e10e1f2d3c79ca32099f1f55c9bc2d43d8}
\newcommand{\leanlink}[1]{%
  \hfill
  \href{\leanrepo/#1}{\normalfont\sffamily\footnotesize Lean 4}%
  \par\nobreak\smallskip}

\newcommand{\Z}{\mathbb Z}
\newcommand{\Q}{\mathbb Q}
\newcommand{\D}{\Q/\Z}
\newcommand{\U}{\operatorname U}
\newcommand{\dlt}{\delta}
\newcommand{\g}{\gamma}
\newcommand{\aug}{\varpi}
\newcommand{\Sym}{\operatorname{Sym}}

\setlength{\abovedisplayskip}{5pt plus 2pt minus 2pt}
\setlength{\belowdisplayskip}{5pt plus 2pt minus 2pt}
\setlength{\abovedisplayshortskip}{3pt plus 2pt}
\setlength{\belowdisplayshortskip}{3pt plus 2pt minus 2pt}

\title{Dimension centrality in Lie algebras and Plotkin's conjecture for metabelian Lie rings}
\author{}
\date{}

\begin{document}

\maketitle

\section{Dimension Centrality}

Let \(k\) be a commutative ring with identity and \(L\) a Lie algebra over
\(k\).  Write \(\U_k(L)\) for its universal enveloping algebra,
\(\iota\colon L\to\U_k(L)\) for the canonical map, and
\(\aug_k(L)=\ker(\U_k(L)\to k)\) for the augmentation ideal.  Put
\[
  \dlt_n(L)=\iota^{-1}\bigl(\aug_k(L)^n\bigr),\qquad
  \g_1(L)=L,\qquad \g_{n+1}(L)=[\g_n(L),L].
\]
All iterated commutators are left-normed.  For an ideal \(A\lhd L\), set
\[
  [A,{}_0L]=A,\qquad [A,{}_{r+1}L]=[[A,{}_rL],L].
\]
No injectivity of \(\iota\), and no PBW theorem, is used in this section.

\begin{center}
\setlength{\fboxsep}{7pt}
\setlength{\fboxrule}{0.7pt}
\fcolorbox{theoremframe}{theoremfill}{%
\begin{minipage}{0.91\linewidth}
\begin{maintheorem}[relative dimension-centrality]\label{thm:relative}
\leanlink{LieRings/DimensionSubring/Centrality.lean\#L29-L37}
For every ideal \(A\lhd L\) and every \(n\geq1\),
\[
  [A,\dlt_n(L)]\subseteq[A,{}_nL].
\]
In particular, \([\dlt_n(L),L]=\g_{n+1}(L)\).
\end{maintheorem}
\end{minipage}}
\end{center}
\vspace{-0.4em}

\begin{proof}
The adjoint action makes \(A\) a right \(\U_k(L)\)-module.  Indeed, on the
tensor algebra of \(L\) define
\[
  a\cdot1=a,\qquad
  a\cdot(x_1\cdots x_r)=[a,x_1,\ldots,x_r].
\]
For \(a\in A\) and \(x,y\in L\), the Jacobi identity gives
\[
 [[a,x],y]-[[a,y],x]-[a,[x,y]]=0,
\]
so the defining relations of \(\U_k(L)\) act trivially.  Since
\(\aug_k(L)^n\) is spanned over \(k\) by monomials of length at least \(n\),
this action satisfies
\[
  A\cdot\aug_k(L)^n=[A,{}_nL].
\]
Indeed, monomials of length \(r\geq n\) act into
\([A,{}_rL]\subseteq[A,{}_nL]\), while those of length exactly \(n\)
give the generators of \([A,{}_nL]\).

If \(a\in A\) and \(x\in\dlt_n(L)\), then
\[
  [a,x]=a\cdot\iota(x)\in A\cdot\aug_k(L)^n=[A,{}_nL],
\]
which proves the relative inclusion.  Taking \(A=L\) and using
antisymmetry gives \([\dlt_n(L),L]\subseteq\g_{n+1}(L)\).  Conversely,
the standard inclusion \(\g_n(L)\subseteq\dlt_n(L)\), obtained by induction
from \(\iota(L)\subseteq\aug_k(L)\) and
\([\aug_k(L)^r,\aug_k(L)]\subseteq\aug_k(L)^{r+1}\), yields
\[
  \g_{n+1}(L)=[\g_n(L),L]\subseteq[\dlt_n(L),L].
\]
\end{proof}

\begin{center}
\setlength{\fboxsep}{7pt}
\setlength{\fboxrule}{0.7pt}
\fcolorbox{corollaryframe}{corollaryfill}{%
\begin{minipage}{0.91\linewidth}
\begin{corollary}\label{cor:mixed-dimension}
\leanlink{LieRings/DimensionSubring/Centrality.lean\#L138-L145}
For all \(m,n\geq1\),
\[
  [\dlt_n(L),\dlt_m(L)]\subseteq\g_{n+m}(L).
\]
\end{corollary}
\end{minipage}}
\end{center}
\vspace{-0.4em}

\begin{proof}
The submodule \(\dlt_n(L)\) is an ideal, so Theorem~\ref{thm:relative}
gives
\[
  [\dlt_n(L),\dlt_m(L)]
  \subseteq[\dlt_n(L),{}_mL].
\]
Iterating the identity
\([\dlt_n(L),L]=\g_{n+1}(L)\) shows that the right-hand side is
\(\g_{n+m}(L)\).
\end{proof}

\section{Metabelian Dimension Vanishing}

We now specialize to \(k=\Z\).  For a Lie ring $L$, let $U(L)$ be its
universal enveloping algebra and
$\aug(L)$ its augmentation ideal.  We use the canonical embedding
$L\hookrightarrow U(L)$ and write
\[
 \delta_s(L)=L\cap\aug(L)^s,
 \qquad
 \gamma_1(L)=L,
 \qquad
 \gamma_{s+1}(L)=[\gamma_s(L),L].
\]
Let $P_jU(L)$ be the subgroup spanned by products of at most $j$
elements of $L$.  Brackets with three or more entries are left-normed,
$\Sym(A)$ means $\Sym_\Z(A)$, and $L$ is metabelian when
$[[L,L],[L,L]]=0$.

\begin{center}
\setlength{\fboxsep}{7pt}
\setlength{\fboxrule}{0.7pt}
\fcolorbox{theoremframe}{theoremfill}{%
\begin{minipage}{0.91\linewidth}
\begin{maintheorem}[metabelian vanishing]\label{thm:nilpotent}
\leanlink{LieRings/DimensionSubring/MetabelianTwoFactor.lean\#L2636-L2641}
If $c\geq2$ is an integer, $L$ is metabelian, and
$\gamma_{c+1}(L)=0$, then
\[
                         \delta_{2c-1}(L)=0.
\]
\end{maintheorem}
\end{minipage}}
\end{center}

Assume first that $L$ is finitely generated as a Lie ring.  Since its
nilpotency class is finite, its additive group is finitely generated.
Put
\[
                         L'=[L,L],
             \qquad L_{\mathrm{ab}}=L/L'.
\]

\medskip\noindent\textbf{1. The free-abelianization cover.}
Choose an exact sequence of abelian groups
\[
 0\longrightarrow K\longrightarrow P
 \overset{\varepsilon}{\longrightarrow}L_{\mathrm{ab}}
 \longrightarrow0,
\]
where $P$ is finitely generated and free, and choose an additive lift
$s:P\to L$ of $\varepsilon$.  Form the pullback
\[
 \widetilde L=L\times_{L_{\mathrm{ab}}}P
 =\{(x,p)\in L\times P:x+L'=\varepsilon(p)\}.
\]
Thus an element of $\widetilde L$ is an element of $L$ together with a
free coordinate for its abelianized class: $\widetilde L$ is obtained
from $L$ by replacing $L_{\mathrm{ab}}$ by $P$.  The first projection
$\pi:\widetilde L\twoheadrightarrow L$ changes neither the derived
ring nor any commutator of weight at least two.  More precisely, after
identifying $K$ with the kernel of $\pi$ and $L'$ with $L'\times0$,
\[
 K\subseteq Z(\widetilde L),\qquad
 \widetilde L'=L',\qquad
 \widetilde L_{\mathrm{ab}}=P,\qquad
 \gamma_j(\widetilde L)\xrightarrow{\ \cong\ }\gamma_j(L)
 \quad(j\geq2).
\]
Via $p\mapsto(s(p),p)$, its additive group is simply $L'\oplus P$;
below we regard elements of $P$ as these chosen lifts.

We shall use the following elementary consequence of metabelianity.
If $d x,d y\in\widetilde L'+Z(\widetilde L)$, then $d[x,y]$ is
central: for every $z\in\widetilde L$,
\[
 [d[x,y],z]=[d x,y,z]=[d x,z,y]=[x,z,d y]=0.
\]
Here central summands may be discarded, and the middle equality says
that the adjoint actions commute on the abelian ring $\widetilde L'$.

Products of additive subgroups below mean the subgroups of the
enveloping algebra generated by the corresponding products.

\begin{lemma}\label{lem:adapted-pbw}
There is an additive retraction
\[
 p_{\leq2}:U(\widetilde L)\longrightarrow P_2U(\widetilde L)
\]
such that
\[
 p_{\leq2}\bigl(KU(\widetilde L)\bigr)
 \subseteq K+K\widetilde L+Z(\widetilde L)\widetilde L.
\]
\end{lemma}

\begin{proof}
Choose a Smith basis of $P$, putting first the basis vectors whose
images in $L_{\mathrm{ab}}$ have finite order, in divisibility order,
and the free ones last.  Put the $L'$-factors first and lift symmetric
monomials in $P$ in the chosen order.  The relations
\[
 ea=ae+[e,a],\qquad ef=fe+[e,f]\quad(f<e),
 \qquad a\in L',
\]
give an additive normal form
\[
 \Sym(L')\otimes_\Z\Sym(P)\xrightarrow{\ \cong\ }U(\widetilde L).
\]
Indeed, each interchange reduces the number of inversions and each
bracket reduces total degree; the only overlaps have length three and
are resolved by the derivation rule and the Jacobi identity.  Since
collection never raises degree, the normal terms of degree at most two
are precisely $P_2U(\widetilde L)$; retaining them defines $p_{\leq2}$.

It remains to prove the displayed inclusion.  Let $e$ be a basis
vector with finite image, let $d$ be the order of $\varepsilon(e)$,
and let $u$ be a homogeneous normal monomial of degree $k$.  The
corresponding basis generator of $K$ has normal expression
\[
                         (0,de)=de+b,\qquad b\in L'.
\]
If $k\leq1$, the retraction leaves $(de+b)u\in K+K\widetilde L$
unchanged.  Suppose $k\geq2$.  The term $bu$ is normal of degree
$k+1$.  When $de$ is collected through $u$, crossing an $L'$-factor
produces no correction, since $d[e,a]=[de+b,a]=0$.  Crossing an
earlier basis vector $f$ produces $d[e,f]$.  The order of
$\varepsilon(f)$ divides $d$.  The Smith relations for $e$ and $f$
therefore give
$de,df\in\widetilde L'+Z(\widetilde L)$, so the observation above
makes $d[e,f]$ central in $\widetilde L$.

The correction-free term has degree $k+1$, while every correction has
degree $k$ and centrality prevents further corrections.  Thus the
projection is zero for $k>2$; for $k=2$ it is a sum of products of a
central element and an element of $\widetilde L$.  This proves the
claim.
\end{proof}

Fix an additive homomorphism $\ell:L\to\D$, and write $\bar x$ for the
image of $x\in\widetilde L$ modulo $Z(\widetilde L)$.  The rule
$(\bar x,\bar y)\mapsto\ell(\pi[x,y])$ is alternating and biadditive.
Ordering the cyclic summands of $\widetilde L/Z(\widetilde L)$ and
taking its upper-triangular part gives a well-defined biadditive map
$H:(\widetilde L/Z(\widetilde L))^2\to\D$ with
\[
 H(\bar x,\bar y)-H(\bar y,\bar x)=\ell(\pi[x,y]).
\]
Define
\[
 \lambda_2:P_2U(\widetilde L)\longrightarrow\D,
 \qquad
 \lambda_2(1)=0,
 \quad \lambda_2(x)=\ell(\pi x),
 \quad \lambda_2(xy)=H(\bar x,\bar y).
\]
This is well defined because its values respect $xy-yx=[x,y]$.
Since $K\subseteq Z(\widetilde L)$,
\[
 \lambda_2(K)=\lambda_2(K\widetilde L)
 =\lambda_2\bigl(Z(\widetilde L)\widetilde L\bigr)=0.
\]
Lemma~\ref{lem:adapted-pbw} therefore gives
\[
 \widetilde\Lambda=\lambda_2\circ p_{\leq2}:
 U(\widetilde L)\longrightarrow\D,
 \qquad
 \widetilde\Lambda\bigl(KU(\widetilde L)\bigr)=0.
\]

\medskip\noindent\textbf{2. The separating trick.}

\begin{lemma}\label{lem:weight-collection}
If $w$ is a word of length $r$ in elements of $\widetilde L$, then
\[
 p_{\leq2}(w)\in\gamma_r(\widetilde L)
 +\sum_{0<a<r}\gamma_a(\widetilde L)\gamma_{r-a}(\widetilde L).
\]
\end{lemma}

\begin{proof}
Collect $w$ by adjacent interchanges $xy=yx+[x,y]$.  Every resulting
term corresponds to a partition of its letters into blocks, a block
of size $a$ giving an element of $\gamma_a(\widetilde L)$.  The
projection retains only terms with one or two blocks.
\end{proof}

\begin{proof}[Proof of Theorem~\ref{thm:nilpotent}]
Lemma~\ref{lem:weight-collection} shows that
\[
                \widetilde\Lambda
                \bigl(\aug(\widetilde L)^{2c-1}\bigr)=0.
\]
Indeed, this power is spanned by words of length $r\geq2c-1$.
Their one-block terms lie in $\gamma_r(\widetilde L)=0$.  In every
two-block term one weight is at least $c$, so the corresponding factor
lies in $\gamma_c(\widetilde L)\subseteq Z(\widetilde L)$ and is
killed by $H$.

The kernel of $U(\widetilde L)\to U(L)$ is the ideal generated by $K$,
which is $KU(\widetilde L)$ because $K$ is central.  Hence
$\widetilde\Lambda$ descends to a homomorphism
\[
 \Lambda:U(L)\longrightarrow\D,
 \qquad
 \Lambda|_L=\ell,
 \qquad
 \Lambda\bigl(\aug(L)^{2c-1}\bigr)=0.
\]
If $0\neq a\in\delta_{2c-1}(L)$, choose $\ell$ with
$\ell(a)\neq0$: a nonzero character of the cyclic group generated by
$a$ extends to $L$ because $\D$ is divisible.  Then
$0=\Lambda(a)=\ell(a)$, a contradiction.

For arbitrary $L$, membership in $L\cap\aug(L)^{2c-1}$ is witnessed in
a finitely generated Lie subring: equivalently, universal enveloping
algebras and augmentation powers commute with filtered colimits.  The
finitely generated case applied there completes the proof.
\end{proof}

\begin{center}
\setlength{\fboxsep}{7pt}
\setlength{\fboxrule}{0.7pt}
\fcolorbox{corollaryframe}{corollaryfill}{%
\begin{minipage}{0.91\linewidth}
\begin{corollary}\label{cor:odd-metabelian}
\leanlink{LieRings/DimensionSubring/MetabelianTwoFactor.lean\#L2912-L2927}
Every metabelian Lie ring $L$ satisfies
$\delta_{2n+1}(L)\subseteq\gamma_{n+2}(L)$ for $n\geq1$.
\end{corollary}
\end{minipage}}
\end{center}
\vspace{-0.4em}

\begin{proof}
Apply Theorem~\ref{thm:nilpotent}, with $c=n+1$, to
$L/\gamma_{n+2}(L)$; the induced map sends $\delta_{2n+1}(L)$ into
$\delta_{2n+1}(L/\gamma_{n+2}(L))=0$.
\end{proof}

\begin{center}
\setlength{\fboxsep}{7pt}
\setlength{\fboxrule}{0.7pt}
\fcolorbox{corollaryframe}{corollaryfill}{%
\begin{minipage}{0.91\linewidth}
\begin{corollary}
\leanlink{LieRings/DimensionSubring/MetabelianTwoFactor.lean\#L2636-L2641}
Every Lie ring $L$ of nilpotency class $n$ satisfies
\[
  \delta_{2n-1}(L)\subseteq L'' \cap Z(L)
\]
\end{corollary}
\end{minipage}}
\end{center}
\vspace{-0.4em}

\begin{proof}
The qutioen $\overline L =L/L''$ is metabelian and has $\gamma_{n+1}(\overline{L}) = 0$, hence $\delta_{2n-1}(\overline L)=0$ by Theorem \ref{thm:nilpotent}. Functoriality yields $\delta_{2n-1}(L)\subseteq L''$. Now appling Theorem \ref{thm:relative} gives $[\delta_{2n-1}(L),L]=\gamma_{2n}(L)\subseteq \gamma_n(L)=0$, hence $\delta_{2n-1}(L)0\subseteq Z(L)$.
\end{proof}
\begin{center}
\setlength{\fboxsep}{7pt}
\setlength{\fboxrule}{0.7pt}
\fcolorbox{corollaryframe}{corollaryfill}{%
\begin{minipage}{0.91\linewidth}
\begin{corollary}
\leanlink{LieRings/DimensionSubring/MetabelianTwoFactor.lean\#L3019-L3031}
Every Lie ring $L$ satisfies
\[
  \delta_5(L)\subseteq\gamma_4(L).
\]
\end{corollary}
\end{minipage}}
\end{center}
\vspace{-0.4em}

\begin{proof}
The quotient $\overline L=L/\gamma_4(L)$ satisfies
$\gamma_4(\overline L)=0$ and is metabelian, since
$\overline L''\subseteq\gamma_4(\overline L)$.  Theorem~\ref{thm:nilpotent},
with $c=3$, gives $\delta_5(\overline L)=0$, and functoriality yields the
claim.
\end{proof}

\begin{center}
\setlength{\fboxsep}{7pt}
\setlength{\fboxrule}{0.7pt}
\fcolorbox{corollaryframe}{corollaryfill}{%
\begin{minipage}{0.91\linewidth}
\begin{corollary}
\leanlink{LieRings/DimensionSubring/MetabelianTwoFactor.lean\#L2934-L2945}
Put
\[
  \delta_\omega(L)=\bigcap_{s\geq1}\delta_s(L),
  \qquad
  \gamma_\omega(L)=\bigcap_{s\geq1}\gamma_s(L).
\]
Then every metabelian Lie ring $L$ satisfies
\[
  \delta_\omega(L)=\gamma_\omega(L).
\]
\end{corollary}
\end{minipage}}
\end{center}
\vspace{-0.4em}

\begin{proof}
The standard inclusions $\gamma_s(L)\subseteq\delta_s(L)$ give
$\gamma_\omega(L)\subseteq\delta_\omega(L)$.  Conversely,
Corollary~\ref{cor:odd-metabelian} gives, for every $n\geq1$,
\[
  \delta_\omega(L)\subseteq\delta_{2n+1}(L)
  \subseteq\gamma_{n+2}(L).
\]
Intersecting over $n$ proves the reverse inclusion.
\end{proof}

\section{Formal Verification}

Theorems~\ref{thm:relative} and~\ref{thm:nilpotent}, together with all the
corollaries above, have been fully verified in Lean~4, with no
\texttt{sorry} declarations and no user-added axioms.  The code is available
in~\cite{IP26}.  We use the definitions of Lie algebras, universal enveloping
algebras, and the lower central series provided by \texttt{mathlib}.  Since
dimension subrings are not currently defined there, we define them by
\[
  \delta_n(L):=\iota^{-1}\bigl(\aug(L)^n\bigr).
\]

To support the formalized proofs, we formalized the
Poincar\'e--Birkhoff--Witt theorem following Higgins's proof~\cite{Hig69}.
In particular, we proved that over $\Z$ the canonical map
$\iota\colon L\hookrightarrow\U_{\Z}(L)$ is injective.  We also formalized
the following results of Bartholdi and Passi~\cite{BP13}:
\[
  \delta_k(L)=\gamma_k(L)\quad(k=2,3),
  \qquad
  2\delta_4(L)\subseteq\gamma_4(L),
\]
as well as $\delta_n(F)=\gamma_n(F)$ for every free Lie ring $F$ and every
$n\geq1$.

A separate part of the development, independent of the present paper,
formalizes Chevalley--Eilenberg homology of Lie algebras through proofs of the
Hopf formula and the five-term exact sequence, following~\cite{IPRZ21}.

\medskip
\begin{center}
\setlength{\fboxsep}{11pt}
\setlength{\fboxrule}{1pt}
\fcolorbox{aiframe}{aifill}{%
\begin{minipage}{0.85\linewidth}
\raggedright
\textbf{AI assistance.}
The OpenAI GPT-5.6 Sol model was used extensively throughout the Lean
formalization.  For the dimension-subring development, the prompts were
submitted in stages: first for the Poincar\'e--Birkhoff--Witt theorem, then
for the Bartholdi--Passi results, next for the corollary
\(\delta_5(L)\subseteq\gamma_4(L)\), and only afterwards for
Theorem~\ref{thm:nilpotent} in its full generality.
\end{minipage}}
\end{center}
\medskip

\begin{thebibliography}{IPRZ21}

\bibitem[BP13]{BP13}
L.~Bartholdi and I.~B.~S.~Passi,
\emph{Lie Dimension Subrings},
\href{https://arxiv.org/abs/1308.2118}{arXiv:1308.2118}, 2013.

\bibitem[Hig69]{Hig69}
P.~J.~Higgins,
\emph{Baer invariants and the Birkhoff--Witt theorem},
J. Algebra \textbf{11} (1969), 469--482;
\href{https://www.cip.ifi.lmu.de/\%7Egrinberg/algebra/higgins-baer.pdf}
{online version}.

\bibitem[IP26]{IP26}
V.~Ionin and T.~Petrov,
\href{https://github.com/TimmPeterson/lie-algebras}
{Lean formalization repository}.

\bibitem[IPRZ21]{IPRZ21}
S.~O.~Ivanov, F.~Pavutnitskiy, V.~Romanovskii, and A.~Zaikovskii,
\emph{On homology of Lie algebras over commutative rings},
J. Algebra \textbf{586} (2021), 99--139;
\href{https://arxiv.org/abs/2010.00369}{arXiv:2010.00369}.

\end{thebibliography}

\end{document}
